% 素粒子物理学 素粒子物理学-2 問題4（014-kosei） 解いた日 2026-10-28
%% kotae
\toi{1}{○}{記述}
スピン，charge，質量

0：クライン・ゴルドン方程式

1/2：ディラック方程式

1：プロカ方程式

角運動量の保存は軌道角運動量とスピンの合計だから。\nuke{空間回転}
\naze{問われている空間回転に触れていない（答案に必要な論点）。空間を回転したときに波動関数がどう変換されるかはスピンで決まるので，スピンが違えば従う波動方程式も異なる，という一言が要る。3つの性質と方程式の対応は正しい。}
\toi{2}{◎}{}
相互作用に対する敏感さ

電：光子，0，電荷

弱：ウィークボソンなど，あり，2種類の弱荷

強：グルーオン，0，カラー荷
\toi{3}{◎}{}
クライン・ゴルドン方程式はスピン0の粒子の方程式だから。

$\pi$中間子，アルファ粒子
%% kizuki
\item スピンは粒子固有の角運動量で，空間回転に対する波動関数の変換のしかたを決める。角運動量の保存（軌道とスピンの合計）は「スピンが何か」の説明で，「なぜ方程式が違うか」には，回転で波動関数がどう変わるかがスピンで違う，という空間回転の側の一言が要る。
\item \mae{第1回 問題4 (3)・10-23} 問題文が「〜に触れて」「〜と比べて」と観点を指定したら，その語を答えに含めたかを最後に確かめる。前回は陽子との比較が，今回は空間回転が抜けた。
